\clearpage
\thispagestyle{empty}
\pdfbookmark[0]{Über ein Mittel, die Beobachtung von Ablenkungen zu erleichtern. Ausgabe und Quellen}{deflection-1839-sources}
\begin{center}\Large\textbf{Über ein Mittel, die Beobachtung von Ablenkungen zu erleichtern}\end{center}

Carl Friedrich Gauss’ Abhandlung erschien in den \textit{Resultaten aus den Beobachtungen des magnetischen Vereins} für 1839 als Artikel II. Die KI-Übertragung aller acht Artikel und der historischen Anmerkung in modernes Deutsch mit KI-erstelltem Satz folgt dem deutschen Text in den \textit{Werken}, Bd.~V (Göttingen, 1867), S.~395--403. Eine Schwingung reicht von einer äußersten Stellung bis zur nächsten. Die Schreibweise der Formeln, die Zahlentafel und die historischen Einheiten folgen der deutschen Ausgabe.

\medskip\noindent
Zur Terminologie der Instrumente siehe den Katalogeintrag \textit{Kommutator} des Deutschen Museums, herausgegeben von Julia Bloemer und Benjamin Mirwald, Nr.~1062, Fassung vom 24. August 2016.

\medskip\noindent
\href{https://gdz.sub.uni-goettingen.de/id/PPN236006339}{Gesammelte deutsche Ausgabe (1867)}\par
\noindent\href{https://digital.deutsches-museum.de/projekte/gruendungssammlung/detail/1062/}{Deutsches Museum, Katalogeintrag zum Kommutator}
\clearpage
